\section{The Same Self, Awaiting Wholeness}

The spiritual soul survives the separation from the body and receives judgment. The Congregation for the Doctrine of the Faith's 1979 letter teaches that consciousness and will remain, so that the human self subsists (no.~3). The one who has loved, sinned, repented, and been forgiven is the one who stands before God. We hope for our own salvation and pray for the people we have known, confident that death has not made them cease to exist or erased the lives for which they will answer.

This continuity gives moral seriousness to the present and solidity to our hope. The acts of charity that appear to vanish into ordinary days belong to a life that endures. So do the choices for which we must ask mercy. God knows the person in the fullness of that history, without the forgetfulness that limits our own account of it. We do not carry ourselves beyond death by the strength of memory. We are held in existence by the God who created us and whose judgment is perfectly true.

Yet a surviving soul is not the whole embodied human being. The blessed truly see God before the general resurrection, while awaiting reunion with their bodies (CCC 1023, citing Benedict XII's \emph{Benedictus Deus}). Their blessedness is real, and humanity's bodily consummation remains future. Thus the Church invokes the saints now and awaits the resurrection of the dead. Our communion with the saints includes hope for what God has still promised to accomplish in them and in us.

Paul's account of the seed and the harvest teaches us to receive continuity together with transformation (1~Cor 15:35--49). The body is sown in weakness and raised in power. What is corruptible is made incorruptible by God. The risen body is our own, yet it will no longer be subject to the conditions under which we now suffer, decline, and die. The language of a spiritual body does not mean that the body has ceased to be bodily. It names embodied life brought fully under the life-giving power of the Spirit, as the risen Christ already reveals.

We can remember this promise when disease, age, or injury has changed a body almost beyond recognition. God will raise the person whose mortal condition we can no longer remedy. Christian reverence therefore remains when beauty, responsiveness, and usefulness have disappeared. The body through which the person has lived deserves care in illness and honor after death. Its decay gives us reason to long for the resurrection God has promised, while the surviving soul remains before him.

The CDF letter warns against arbitrary pictures of the interval between death and resurrection (nos.~2--5 and concluding doctrinal paragraphs). It affirms the truths on which the Church's prayer depends: the self endures, judgment is real, the saints live in God, the departed can be helped by intercession, and bodily resurrection is still to come. These are definite grounds for hope even when we cannot answer further questions about the manner of life after death.

Beside a bed or grave we can commend the person by name, ask mercy, honor the body, and await the fulfillment of God's promise. We need not know the details of the departed person's condition in order to pray. The Church continues to intercede for the dead because Christ makes that prayer effective; we can participate in it even when grief leaves us little else to say.

